\chapter{Gérer un portefeuille pendant un drawdown}\label{ch:fm:managing-a-book-through-drawdown}

Deux gérants de portefeuille sont tous deux en recul de huit pour cent par rapport à leur sommet de mars. L'un fait tourner une stratégie qui a connu de mauvais mois~; chez le second,
l'avantage s'est éteint il y a un an, et personne ne s'en est aperçu. Le jour où chacun atteint pour la première fois les huit pour cent, la meilleure lecture statistique de leur P\&L donne
au premier une probabilité de 24~\% d'avoir décroché et au second de 39~\%~: les chiffres ne diffèrent guère. Pour distinguer une stratégie dont le ratio de Sharpe est passé de 1
à 0 de celle d'une stratégie simplement malchanceuse, à raison d'une fausse alarme par décennie, il faut en moyenne plus de quinze ans de données. Le responsable du desk
doit trancher en mars.

\section{Ce que dit un drawdown sur un ratio de Sharpe}

Un drawdown (livre~2, chapitre~29) est la baisse du P\&L cumulé depuis son plus haut courant~; son maximum sur une période et sa durée sont
les mesures du livre~7 (chapitre~22). Ce qui compte pour la gestion, c'est l'ampleur du drawdown qu'engendre une stratégie \emph{saine}. Pour un P\&L quotidien avec un
ratio de Sharpe de 1 et une volatilité de 10~\% par an sur le capital, le drawdown maximal sur dix ans a une médiane de 16.8~\% du capital et un
90\up{e}~centile de 25.4~\% (2\,000 trajectoires simulées)~: une stratégie saine de ce type passera quelques années à plus d'une année de volatilité sous
son sommet. Magdon-Ismail, Atiya, Pratap et Abu-Mostafa donnent la distribution du drawdown maximal d'un mouvement brownien avec dérive~; la
simulation reproduit sa leçon~: des drawdowns d'une à deux années de volatilité sont le régime normal d'une stratégie de Sharpe~1.

\begin{remark}[Unités]\label{rem:fm:managing-a-book-through-drawdown:units}
Un drawdown ne veut rien dire tant qu'on ne le rapporte pas à la volatilité de la stratégie. Huit pour cent valent 0.8 année de volatilité à 10~\%, et 0.27 à 30~\%.
Toutes les règles de ce chapitre sont énoncées pour une stratégie de volatilité de 10~\%~; pour une autre, il faut proportionner les seuils au rapport des volatilités (le
livre~8, chapitre~28, est arrivé à la même conclusion en partant de l'autre côté).
\end{remark}

\section{Règles de stop-loss et échelles de réduction du risque}

\begin{definition}[Règle de stop-loss, échelle de réduction du risque, stop temporel]\label{def:fm:managing-a-book-through-drawdown:rules}
Une \emph{règle de stop-loss}\index{règle de stop-loss} (stop-loss rule) arrête une stratégie lorsque son drawdown par rapport à son plus haut courant dépasse un seuil. Une \emph{échelle de
réduction du risque}\index{échelle de réduction du risque} (de-risking ladder) réduit la taille d'une stratégie par paliers à mesure que son drawdown s'approfondit (division par deux au premier seuil, arrêt au second)
au lieu de l'arrêter d'un coup. Un \emph{stop temporel}\index{stop temporel} (time stop) arrête une stratégie qui n'a inscrit aucun nouveau plus haut depuis une durée donnée,
quelle que soit la profondeur de son drawdown.
\end{definition}

Une \omterm{def:fm:managing-a-book-through-drawdown:rules}{règle de stop-loss} commet deux erreurs, et la limite de drawdown du livre~8, chapitre~28, en avait déjà montré la forme~: elle arrête des stratégies en production qui traversent une passe de
malchance, et laisse courir les stratégies mortes. L'échelle est une version à deux paliers de la règle de Grossman et Zhou, où la position diminue en douceur à mesure que le
drawdown approche d'un plancher~; le \omterm{def:fm:managing-a-book-through-drawdown:rules}{stop temporel} surveille la durée plutôt que la profondeur, car une stratégie morte, de ratio de Sharpe nul, dérive
latéralement plutôt que vers le bas.

\omcode{code/firm/ddrules/firm_ddrules.py}{40}{55}{La probabilité a posteriori qu'une stratégie ait décroché~: un filtre progressif sur une chaîne à deux états,
active et décrochée, avec un taux de décrochage quotidien.}\label{lst:fm:managing-a-book-through-drawdown:posterior}

\section{Malchance ou décrochage~: preuves séquentielles}

Une règle de drawdown ne regarde qu'une seule statistique de la trajectoire du P\&L. L'ensemble de la trajectoire en contient davantage~: le rendement de chaque jour apporte un élément de preuve sur la question de savoir si la stratégie
est encore celle qui a été backtestée.

\begin{proposition}[La probabilité a posteriori qu'une stratégie ait décroché]\label{prop:fm:managing-a-book-through-drawdown:posterior}
Soit une stratégie active ou décrochée~; active, son rendement quotidien suit $\mathcal N(\mu_1,\sigma_d^2)$~; décrochée, $\mathcal N(\mu_0,\sigma_d^2)$~; elle décroche
un jour quelconque avec probabilité $h$ si elle ne l'est pas encore, et ne se rétablit jamais. Avec $p_{t-1}$ la probabilité qu'elle soit décrochée sachant les rendements jusqu'au jour
$t-1$, la probabilité, sachant le jour $t$ et son rendement $x_t$, est
\[
p_t=\frac{\tilde p_t\,\varphi\!\left(\frac{x_t-\mu_0}{\sigma_d}\right)}{\tilde p_t\,\varphi\!\left(\frac{x_t-\mu_0}{\sigma_d}\right)+(1-\tilde p_t)\,\varphi\!\left(\frac{x_t-\mu_1}{\sigma_d}\right)},
\qquad \tilde p_t=p_{t-1}+(1-p_{t-1})\,h .
\]
\end{proposition}

\begin{proof}
L'étape de prédiction~: la stratégie est décrochée au jour $t$ si elle l'était auparavant ou si elle décroche aujourd'hui, $\tilde p_t=p_{t-1}+(1-p_{t-1})h$. La mise à jour est
la règle de Bayes avec les deux vraisemblances gaussiennes~; le facteur commun $1/\sigma_d$ s'annule.
\end{proof}

\begin{proposition}[Combien de temps il faut pour savoir]\label{prop:fm:managing-a-book-through-drawdown:lorden}
Pour la détection d'un changement de la loi active à la loi décrochée, avec une moyenne de $A$ jours entre fausses alarmes, aucun détecteur n'a, lorsque
$A\to\infty$, un délai moyen dans le pire des cas inférieur à $\ln A/K$, et le détecteur CUSUM l'atteint, où $K$ est la divergence de Kullback--Leibler
par jour. Pour des rendements quotidiens de ratios de Sharpe annuels $S$ avant et $S_0$ après le décrochage, $K=(S-S_0)^2/(2\times252)$, le délai vaut
\[
\frac{2\,\ln A}{(S-S_0)^2}\ \text{années}.
\]
\end{proposition}

\admitted

La borne est le résultat de Lorden (1971), cité dans \texttt{omsources}~; le test CUSUM lui-même relève du livre~7 (chapitre~13). Son calcul illustre le
fait central du chapitre. Un ratio de Sharpe qui passe de 1 à 0, détecté à raison d'une fausse alarme par décennie ($A=2\,520$ jours), demande
$2\ln2\,520=15.7$ années pour être détecté en moyenne~; à raison d'une fausse alarme tous les cinq ans, 14.3. Une baisse à $-1$ demande 3.9 années et une baisse à $-2$, 1.7
(la \cref{fig:fm:managing-a-book-through-drawdown:lorden}). Un avantage qui disparaît discrètement ne peut se distinguer de la malchance sur son P\&L dans aucun
délai utile~; seul un avantage qui se transforme en perte le peut.

\begin{omfigure}
\begin{tikzpicture}
\begin{axis}[width=10.5cm,height=5.2cm,xlabel={\small Ratio de Sharpe après le décrochage $S_0$ (avant~: 1)},ylabel={\small années avant détection},
  xmin=-2,xmax=0.5,ymin=0,ymax=65,tick label style={font=\footnotesize},grid=major,grid style={gray!20},table/col sep=comma]
\addplot[omDef,thick] table[x=sr_dead,y=years]{figdata/desk/08-managing-a-book-through-drawdown/lorden.csv};
\draw[omThm,dashed] (axis cs:0,0) -- (axis cs:0,65);
 \node[font=\scriptsize,anchor=west,fill=white,inner sep=1pt] at (axis cs:-1.9,50) {une fausse alarme par décennie};
\end{axis}
\end{tikzpicture}

\omcaption{Le délai moyen le plus court avec lequel un détecteur quelconque peut déceler que le ratio de Sharpe d'une stratégie est tombé de 1 à $S_0$, à raison d'une fausse
alarme par décennie, d'après la \cref{prop:fm:managing-a-book-through-drawdown:lorden}~: 15.7 années pour une baisse jusqu'à zéro (trait en pointillés), 3.9 pour une baisse jusqu'à
$-1$. Données~: \texttt{fm\_drawdown.lorden}.}\label{fig:fm:managing-a-book-through-drawdown:lorden}
\end{omfigure}

La probabilité a posteriori de la \cref{prop:fm:managing-a-book-through-drawdown:posterior}, calculée sur deux trajectoires d'exemple, montre ce que voit le responsable du desk
(la \cref{fig:fm:managing-a-book-through-drawdown:paths}). La probabilité a posteriori de la stratégie active oscille, atteignant 0.80 la quatrième année~; celle de la
stratégie décrochée croît lentement après son décrochage, au début de la quatrième année, et ne dépasse 0.9 qu'à sa cinquième année.

\begin{omfigure}
\begin{tikzpicture}
\begin{axis}[width=10.5cm,height=4.4cm,ylabel={\small P\&L cumulé (\%)},xmin=0,xmax=10,ymin=-20,ymax=100,
  tick label style={font=\footnotesize},grid=major,grid style={gray!20},table/col sep=comma,xticklabels={},
  legend cell align=left,legend style={font=\scriptsize,at={(0.02,0.97)},anchor=north west}]
\addplot[omDef,thick] table[x=year,y=eq_alive]{figdata/desk/08-managing-a-book-through-drawdown/examples.csv};
\addplot[omThm,thick,dashed] table[x=year,y=eq_broken]{figdata/desk/08-managing-a-book-through-drawdown/examples.csv};
\draw[gray] (axis cs:3,-20) -- (axis cs:3,100);
\legend{active,{décroche à l'année 3}}
\end{axis}
\begin{axis}[width=10.5cm,height=4.4cm,yshift=-3.6cm,xlabel={\small année},ylabel={\small P(décroché)},xmin=0,xmax=10,ymin=0,ymax=1.05,
  tick label style={font=\footnotesize},grid=major,grid style={gray!20},table/col sep=comma]
\addplot[omDef,thick] table[x=year,y=post_alive]{figdata/desk/08-managing-a-book-through-drawdown/examples.csv};
\addplot[omThm,thick,dashed] table[x=year,y=post_broken]{figdata/desk/08-managing-a-book-through-drawdown/examples.csv};
\draw[gray] (axis cs:3,0) -- (axis cs:3,1.05);
\draw[omProp,dotted,thick] (axis cs:0,0.9) -- (axis cs:10,0.9);
\end{axis}
\end{tikzpicture}

\omcaption{Deux stratégies de ratio de Sharpe 1 et de volatilité 10~\%~; le ratio de Sharpe de la seconde tombe à 0 au début de l'année 4 (ligne
verticale). En haut~: P\&L cumulé. En bas~: la probabilité a posteriori que chacune ait décroché (la \cref{prop:fm:managing-a-book-through-drawdown:posterior},
avec un taux de hasard d'une rupture tous les cinq ans), le seuil de 0.9 en pointillés. Données~: \texttt{fm\_drawdown.examples}.}\label{fig:fm:managing-a-book-through-drawdown:paths}
\end{omfigure}

\section{Comparer les règles}

L'évaluation du chapitre porte sur 2\,000 stratégies actives et sur 2\,000 stratégies qui décrochent au début de l'année 4 et ramènent leur ratio de Sharpe à 0, pendant dix ans, sous chaque
règle, tous les arrêts étant définitifs. Elle rapporte les faux arrêts pour cent années-stratégie actives, la part des stratégies décrochées arrêtées après leur décrochage,
et le délai médian, ainsi que la valeur conservée par stratégie sur dix ans~: le P\&L diminué d'une \omterm{def:fm:limits-and-capital-allocation:raroc}{charge en capital} de 2~\% du capital pour chaque année exploitée à pleine
taille (un seuil de rendement de 15~\% sur un capital égal à 13~\% de la taille de la stratégie), en pourcentage du capital.

\begin{center}
\small
\begin{tabular}{@{}>{\raggedright\arraybackslash}p{2.6cm}>{\centering\arraybackslash}p{1.75cm}>{\centering\arraybackslash}p{1.75cm}>{\centering\arraybackslash}p{1.75cm}>{\centering\arraybackslash}p{1.75cm}>{\centering\arraybackslash}p{1.75cm}>{\centering\arraybackslash}p{1.75cm}@{}}
\toprule
règle & faux arrêts & décrochées & médiane & valeur, & valeur, & valeur,\\
 & pour 100 années & détectées & délai (années) & actives & décrochées & moitié chacune\\
\midrule
aucune & -- & -- & -- & 79.0 & 11.1 & 45.1\\
stop-loss 10~\% & 9.79 & 34~\% & 0.72 & 21.7 & 15.0 & 18.3\\
échelle 7.5 / 12.5~\% & 2.09 & 60~\% & 2.79 & 39.5 & 12.8 & 26.2\\
\omterm{def:fm:managing-a-book-through-drawdown:rules}{stop temporel} 18 mois & 6.19 & 80~\% & 2.00 & 54.6 & 19.5 & 37.1\\
a posteriori au-dessus de 0.9 & 2.84 & 95~\% & 2.68 & 71.4 & 18.2 & 44.8\\
\bottomrule
\end{tabular}
\end{center}

\begin{omfigure}
\begin{tikzpicture}
\begin{axis}[width=11cm,height=5.2cm,ybar,bar width=8pt,ymin=0,ymax=85,ylabel={\small valeur conservée (\% du capital)},
  xtick={0,1,2,3,4},xticklabels={aucune,{stop-loss 10~\%},{échelle 7.5 / 12.5~\%},{stop temporel 18 mois},{a posteriori $>0.9$}},
  xticklabel style={font=\scriptsize,text width=1.8cm,align=center},tick label style={font=\footnotesize},table/col sep=comma,
  enlarge x limits=0.12,grid=major,grid style={gray!20},legend cell align=left,
  legend style={font=\scriptsize,at={(0.5,-0.32)},anchor=north,/tikz/every even column/.append style={column sep=8pt}},legend columns=3]
\addplot[fill=omDef!55,draw=omDef] table[x=k,y=value_alive]{figdata/desk/08-managing-a-book-through-drawdown/rules.csv};
\addplot[fill=omThm!50,draw=omThm] table[x=k,y=value_dead]{figdata/desk/08-managing-a-book-through-drawdown/rules.csv};
\addplot[fill=gray!40,draw=gray] table[x=k,y=value_mix]{figdata/desk/08-managing-a-book-through-drawdown/rules.csv};
\legend{stratégies actives,stratégies décrochées,moitié de chaque}
\end{axis}
\end{tikzpicture}

\omcaption{Valeur conservée par stratégie sur dix ans (P\&L diminué d'une \omterm{def:fm:limits-and-capital-allocation:raroc}{charge en capital} de 2~\%-par-an pendant l'exploitation, en pourcentage du capital) pour chaque règle,
pour les stratégies qui restent actives, les stratégies qui décrochent à l'année 3 et ramènent leur ratio de Sharpe à 0, et une population de moitié chacune. 2\,000 trajectoires de chaque type~; les arrêts
sont définitifs. Données~: \texttt{fm\_drawdown.table}.}\label{fig:fm:managing-a-book-through-drawdown:rules}
\end{omfigure}

Le stop-loss à 10~\% est de loin la pire des règles. Il arrête 98~\% des stratégies saines à un moment donné sur dix ans, et les deux tiers des stratégies décrochées
\emph{avant} leur décrochage~; il conserve 18.3~\% du capital par stratégie, contre 45.1~\% sans aucune règle. La règle a posteriori arrête 28~\% des stratégies saines
sur dix ans, détecte 95~\% des stratégies décrochées environ deux ans et demi après le décrochage, et conserve 44.8~\% du capital, autant que
de ne rien faire. Arrêter une stratégie décrochée dont le ratio de Sharpe est nul ne fait économiser aucun P\&L en espérance~; elle économise le capital immobilisé et le risque ajouté,
ce qui explique que la \omterm{def:fm:limits-and-capital-allocation:raroc}{charge en capital} figure dans la valeur. Avec une charge nulle, aucune règle ne fait mieux que de ne rien faire~; pour une stratégie qui devient perdante, chaque
règle la détecte plus tôt, et la règle a posteriori est nettement la meilleure (exercice 7).

\begin{method}[Revue d'une stratégie en drawdown]\label{met:fm:managing-a-book-through-drawdown:review}
\begin{enumerate}
\item Exprimer le drawdown et sa durée en unités de la volatilité de la stratégie~; les comparer à la distribution d'une stratégie saine.
\item Mettre à jour la probabilité a posteriori que la stratégie ait décroché à partir de tout son historique de P\&L, avec un taux de hasard calé sur l'expérience de la firme en matière de durée
pendant laquelle les stratégies restent actives (livre~11, chapitre~28).
\item Chercher des indices qui ne sont pas du P\&L~: taux de capture et taux d'exécution, mesures d'encombrement (crowding) (livre~7, chapitre~28), un changement dans la structure du
marché. Ces indices séparent bien plus vite la malchance du décrochage que ne le fait le P\&L.
\item Réduire la taille progressivement à mesure que les indices s'accumulent~; n'arrêter que sur la probabilité a posteriori ou sur des indices autres que le P\&L~; jamais sur la seule profondeur.
\end{enumerate}
\end{method}

\section{Redémarrer, réévaluer, retirer}

\begin{definition}[Règle de redémarrage]\label{def:fm:managing-a-book-through-drawdown:restart}
Une \emph{règle de redémarrage}\index{règle de redémarrage} précise quand et à quelle taille une stratégie arrêtée peut de nouveau trader~: après une période fixe, après que son
P\&L papier se redresse, ou après qu'un réexamen de sa thèse la juge intacte.
\end{definition}

Un arrêt est une décision prise en situation d'incertitude~; une \omterm{def:fm:managing-a-book-through-drawdown:restart}{règle de redémarrage} est ce qui la rend réversible. Comme une stratégie saine est arrêtée par toute règle qui
détecte les stratégies décrochées dans un délai utile, le coût des faux arrêts se paie dans les mois pendant lesquels la stratégie est à plat. Une \omterm{def:fm:managing-a-book-through-drawdown:restart}{règle de redémarrage} qui fait du paper trading sur une stratégie arrêtée
et la remet en marche à taille réduite quand son P\&L papier se redresse récupère une bonne part de ce coût~; un redémarrage qui intervient automatiquement
à la nouvelle année, comme dans le modèle du livre~8, chapitre~28, ne le récupère que pour les bonnes stratégies et relance aussi les stratégies décrochées.
\emph{La réévaluation} d'une stratégie consiste à l'examiner comme si elle était nouvelle~: sa thèse, ses indices, sa capacité et son coût du capital. Une
stratégie qui échoue à cet examen est retirée même si son P\&L ne l'a pas encore dit, seule façon de battre les délais de
la \cref{prop:fm:managing-a-book-through-drawdown:lorden}.

\begin{remark}[Les acteurs en drawdown]\label{rem:fm:managing-a-book-through-drawdown:people}
Un gérant de portefeuille en drawdown est soumis à des règles publiques et asymétriques~: un arrêt met fin à son poste, un redressement ne rétablit la rémunération qu'au-dessus du
high-water mark (chapitre~3). Les deux poussent à prendre davantage de risque pour se redresser vite, ou moins pour éviter l'arrêt. Une échelle qui réduit la taille avant l'arrêt
et une \omterm{def:fm:managing-a-book-through-drawdown:restart}{règle de redémarrage} qui ne met pas fin à une carrière dès le premier arrêt réduisent l'incitation à miser sur une résurrection~; il en va de même d'une firme qui est
connue pour examiner les thèses, et pas seulement le P\&L.
\end{remark}

\section{Tutoriel~: deux mille stratégies et quatre règles}\label{tut:fm:managing-a-book-through-drawdown:tut}

\begin{tutorial}
\textbf{Objectif.} Mesurer ce que coûte chaque règle de drawdown à une stratégie saine et ce qu'elle fait économiser sur une stratégie décrochée. \textbf{État final~:} le tableau des règles et
la \cref{fig:fm:managing-a-book-through-drawdown:rules}.
\begin{tutsteps}
\item \textbf{Trajectoires.} \texttt{fm\_drawdown.population()} tire 2\,000 stratégies actives et 2\,000 stratégies qui décrochent avec \texttt{firm.ddrules.paths}.
\item \textbf{A posteriori.} \texttt{firm.ddrules.posterior} applique le filtre du \cref{lst:fm:managing-a-book-through-drawdown:posterior} à chaque
trajectoire.
\item \textbf{Règles.} \texttt{fm\_drawdown.table()} applique les quatre règles avec \texttt{apply} et les note avec \texttt{evaluate}.
\item \textbf{La borne.} \texttt{fm\_drawdown.lorden(10, S0)} évalue la \cref{prop:fm:managing-a-book-through-drawdown:lorden}.
\end{tutsteps}
\end{tutorial}

\textbf{Pour aller plus loin.} Faire décrocher les stratégies jusqu'à un ratio de Sharpe de $-1$ (exercice 7)~; fixer le taux de hasard dix fois trop haut et observer comment la
règle a posteriori arrête des stratégies saines.

\section{Construction~: règles de drawdown}\label{bld:fm:managing-a-book-through-drawdown:ddrules}

\begin{build}
\bfield{Rôle} Règles de gestion des stratégies en drawdown, et une mesure honnête de ce que chacune coûte et fait économiser.
\bfield{Interface} \texttt{firm.ddrules}~: \texttt{paths(n, days, sr, vol, rng, break\_day, sr\_dead)}~;
\texttt{posterior(returns, sr, sr\_dead, vol, hazard)}~; \texttt{Rule(kind, a, b)} avec pour types \texttt{stop}, \texttt{ladder}, \texttt{time}, \texttt{posterior}~; \texttt{apply}~;
\texttt{evaluate(rule, alive, dead, break\_day, post\_alive, post\_dead, charge)}~; \texttt{cusum\_delay(sr, sr\_dead, arl\_days)}.
\bfield{Règles} La taille d'un jour se décide sur le P\&L jusqu'à la veille~; les arrêts sont définitifs dans l'évaluateur~; la valeur déduit la \omterm{def:fm:limits-and-capital-allocation:raroc}{charge en capital}
seulement pendant que la stratégie tourne.
\bfield{Tests d'acceptation} \texttt{code/firm/ddrules/tests/}~: chaque règle sur une trajectoire faite à la main (le jour de l'arrêt, la division par deux, le \omterm{def:fm:managing-a-book-through-drawdown:rules}{stop temporel})~; la
probabilité a posteriori sépare les populations actives des populations décrochées~; la forme de l'évaluateur~; la formule du délai.
\bfield{Pour aller plus loin} \omterm{def:fm:managing-a-book-through-drawdown:restart}{Règles de redémarrage} dans l'évaluateur~; une distribution a posteriori sur le ratio de Sharpe après décrochage au lieu d'un $S_0$ fixe~; des indices tirés des taux d'exécution
et de capture comme un second flux d'observations.
\end{build}

\begin{omsources}
\item M.~Magdon-Ismail, A.~F.~Atiya, A.~Pratap et Y.~S.~Abu-Mostafa, ``On the maximum drawdown of a Brownian motion'', \emph{Journal of
Applied Probability} 41(1), 2004.
\item G.~Lorden, ``Procedures for reacting to a change in distribution'', \emph{Annals of Mathematical Statistics} 42(6), 1971.
\item R.~P.~Adams et D.~J.~C.~MacKay, ``Bayesian online changepoint detection'', arXiv:0710.3742, 2007.
\end{omsources}

\section{Exercices}

\begin{exercise}[$\star$]\label{exo:fm:managing-a-book-through-drawdown:1}
Exprimer un drawdown de 8~\% en années de volatilité pour des stratégies de volatilités de 10~\% et 30~\%.
\end{exercise}

\begin{exercise}[$\star$]\label{exo:fm:managing-a-book-through-drawdown:2}
Utiliser la \cref{prop:fm:managing-a-book-through-drawdown:lorden} pour calculer le délai de détection d'une baisse du ratio de Sharpe de 1 à 0 avec une fausse alarme
par dix ans, et de 1 à $-1$.
\end{exercise}

\begin{exercise}[$\star$]\label{exo:fm:managing-a-book-through-drawdown:3}
D'après le tableau des règles, quelle part des stratégies saines le stop-loss à 10~\% arrête-t-il sur dix ans, et quelle part des stratégies décrochées arrête-t-il
avant leur décrochage~?
\end{exercise}

\begin{exercise}[$\star\star$]\label{exo:fm:managing-a-book-through-drawdown:4}
La probabilité a posteriori qu'une stratégie ait décroché vaut 0.30 et le hasard quotidien vaut $1/1\,260$. Son rendement du lendemain est exactement la moyenne des stratégies décrochées. La probabilité
a posteriori est-elle plus haute ou plus basse après ce jour, et pourquoi un seul jour la déplace-t-il si peu~?
\end{exercise}

\begin{exercise}[$\star\star$]\label{exo:fm:managing-a-book-through-drawdown:5}
Pourquoi arrêter une stratégie décrochée dont le ratio de Sharpe est nul n'ajoute-t-il aucun P\&L espéré, et que permet-il d'économiser~?
\end{exercise}

\begin{exercise}[$\star\star$]\label{exo:fm:managing-a-book-through-drawdown:6}
Une stratégie a un ratio de Sharpe de 1 et une volatilité de 10~\%. Son drawdown maximal sur dix ans a une médiane de 16.8~\%. Un drawdown de 12~\% dans sa
sixième année est-il une preuve qu'elle a décroché~?
\end{exercise}

\begin{exercise}[$\star\star\star$]\label{exo:fm:managing-a-book-through-drawdown:7}
\emph{Programmation.} Relancer l'évaluation en faisant chuter les stratégies décrochées à un ratio de Sharpe de $-1$ (fixer \texttt{SR\_DEAD} et l'alternative de la
probabilité a posteriori à $-1$). Que deviennent les délais et la valeur de chaque règle~?
\end{exercise}

\begin{exercise}[$\star\star\star$]\label{exo:fm:managing-a-book-through-drawdown:8}
\emph{Trouver la faille.} «~Nous avons testé notre stop-loss à 10~\% sur l'historique de nos stratégies, et il aurait évité chaque grosse perte. Adoptons-le.~»
\end{exercise}

\section{Problème~: Huit points de baisse en mars}

\begin{problem}[{Problème du week-end --- Huit points de baisse en mars}]\label{pb:fm:managing-a-book-through-drawdown:1}
Deux gérants de portefeuille perdent chacun 8~\% depuis leur sommet. Le responsable du desk doit décider quoi faire de chacun, et quelle règle adopter pour la prochaine fois.

\medskip\noindent\textbf{Partie I --- Le drawdown.}
\begin{enumerate}
\item Définir une \omterm{def:fm:managing-a-book-through-drawdown:rules}{règle de stop-loss}, une \omterm{def:fm:managing-a-book-through-drawdown:rules}{échelle de réduction du risque} et un \omterm{def:fm:managing-a-book-through-drawdown:rules}{stop temporel}.
\item Donner la médiane et le 90\up{e}~centile du drawdown maximal sur dix ans d'une stratégie saine de Sharpe~1 à 10~\% de volatilité.
\item Exprimer 8~\% en unités de volatilité pour les deux gérants, à 10~\% et 20~\% de volatilité.
\item Pourquoi les seuils doivent-ils s'exprimer en unités de volatilité~?
\end{enumerate}

\medskip\noindent\textbf{Partie II --- Les indices.}
\begin{enumerate}[resume]
\item Énoncer et démontrer la \cref{prop:fm:managing-a-book-through-drawdown:posterior}.
\item Donner la probabilité a posteriori moyenne le jour où une stratégie saine et une stratégie décrochée atteignent pour la première fois un drawdown de 8~\%.
\item Énoncer la \cref{prop:fm:managing-a-book-through-drawdown:lorden} et calculer le délai pour une baisse de 1 à 0 et à $-1$ avec une fausse alarme par
dix ans.
\item Quels types d'indices distinguent la malchance du décrochage plus vite que le P\&L~?
\end{enumerate}

\medskip\noindent\textbf{Partie III --- Les règles.}
\begin{enumerate}[resume]
\item Donner, pour chaque règle, le nombre de fausses alertes pour cent années-stratégies et la part des stratégies décrochées arrêtées après leur décrochage.
\item Donner, pour chaque règle, le délai médian de détection.
\item Donner, pour chaque règle, la valeur conservée pour les populations saines, décrochées et mixtes, face à l'absence de règle.
\item Pourquoi le stop-loss est-il la pire des règles~?
\item Pourquoi la règle fondée sur la probabilité a posteriori conserve-t-elle à peu près autant que l'absence de règle, et quand conserverait-elle davantage~?
\end{enumerate}

\medskip\noindent\textbf{Partie IV --- La décision.}
\begin{enumerate}[resume]
\item Définir une \omterm{def:fm:managing-a-book-through-drawdown:restart}{règle de redémarrage} et en décrire une bonne.
\item Qu'est-ce que la re-souscription, et pourquoi peut-elle battre les délais du P\&L~?
\item Que font les règles aux incitations d'un gérant de portefeuille en drawdown~?
\item Que feriez-vous de chacun des deux gérants lundi~?
\item Quelle règle adopteriez-vous, avec quels seuils, et pourquoi~?
\item Énoncer le \emph{résultat nommé}~: la probabilité a posteriori le jour du drawdown de 8~\% pour une stratégie saine et une stratégie décrochée, et le délai de détection
d'une baisse de 1 à 0 avec une fausse alarme par dix ans.
\item En deux phrases, expliquer aux gérants comment ils seront jugés.
\end{enumerate}
\end{problem}

\section{Questions d'entretien}

\begin{interviewq}[$\star$ \iqroles{trader}]\label{iq:fm:managing-a-book-through-drawdown:1}
Votre stratégie a un ratio de Sharpe de 1 et une volatilité de 10~\%, et elle se situe 15~\% sous son sommet. Devez-vous vous inquiéter~?
\end{interviewq}

\begin{interviewq}[$\star$ \iqroles{risque}]\label{iq:fm:managing-a-book-through-drawdown:2}
Quelle est la différence entre un stop-loss et une \omterm{def:fm:managing-a-book-through-drawdown:rules}{échelle de réduction du risque}, et pourquoi une firme préférerait-elle l'échelle~?
\end{interviewq}

\begin{interviewq}[$\star\star$ \iqroles{chercheur}]\label{iq:fm:managing-a-book-through-drawdown:3}
Combien de temps faudrait-il pour détecter, à partir du P\&L quotidien, qu'un ratio de Sharpe est passé de 1 à 0~? Dérivez l'ordre de grandeur.
\end{interviewq}

\begin{interviewq}[$\star\star$ \iqroles{chercheur, risque}]\label{iq:fm:managing-a-book-through-drawdown:4}
Écrivez la mise à jour de la probabilité qu'une stratégie ait décroché, étant donné un jour de rendements supplémentaire.
\end{interviewq}

\begin{interviewq}[$\star\star$ \iqroles{trader, risque}]\label{iq:fm:managing-a-book-through-drawdown:5}
Un backtest montre qu'un stop-loss aurait évité chaque grosse perte. Que vérifieriez-vous~?
\end{interviewq}

\begin{interviewq}[$\star\star\star$ \iqroles{chercheur}]\label{iq:fm:managing-a-book-through-drawdown:6}
Votre P\&L ne peut pas vous dire à temps si un edge a disparu. Que surveilleriez-vous d'autre, et pourquoi est-ce plus rapide~?
\end{interviewq}
